% =====================================================================
%  Treść — MorphBPE: morfologicznie informowany BPE dla polszczyzny
% =====================================================================

\begin{abstract}
\noindent
Prezentujemy MorphBPE --- wariant algorytmu Byte-Pair Encoding, który
wymusza poszanowanie granic morfemów podczas scalania par tokenów.
Granice sufiksów pozyskiwane są z~analizatora morfologicznego Morfeusz2;
scalenia przekraczające granicę sufiksu są karane mnożnikiem częstotliwości.
Porównujemy MorphBPE (wariant \texttt{strict\_suffix}) ze standardowym BPE
na korpusie 10\,000 przemówień sejmowych (split 80/20, seed=42)
przy rozmiarach słownika 500 i~1000.
Ewaluacja intrinsic (Fertility, Morphological Edit Distance) z~bootstrap
CI (1000 resamples, 95\%) pokazuje, że MorphBPE redukuje MorphED
o~3{,}6\%--6{,}8\% bez kosztu fertility.
Ewaluacja extrinsic z~dwoma architekturami --- LSTM (1 warstwa)
i~Transformer (2 warstwy, 4 głowice) --- potwierdza,
że MorphBPE obniża cross-entropy o~1{,}3\%--2{,}7\%
i~perplexity o~4{,}0\%--9{,}3\% względem BPE przy tym samym rozmiarze słownika.
Analiza błędów per-suffix pokazuje, że MorphBPE zyskuje
na sufiksach złożonych (\texttt{-ach}: $-11{,}5\%$),
ale traci na krótkich (\texttt{-ych}: $+42{,}1\%$).
Rozszerzone badanie siły kary $\alpha \in [0{,}0;\,3{,}0]$
pokazuje, że $\alpha = 2{,}0$ (tzw.\ over-penalization) daje
lepsze wyniki niż standardowy zakres $\alpha \le 1{,}0$:
CE spada o~dodatkowe $1{,}6\%$--$3{,}1\%$ względem $\alpha = 1{,}0$.
Skalowanie do rozmiarów produkcyjnych (vocab 8\,000--32\,000)
potwierdza trend: redukcja MorphED rośnie do $21{,}0\%$,
a~over-penalization pozostaje optymalne na wszystkich rozmiarach.
\end{abstract}

% ---------------------------------------------------------------
\section{Przegląd literatury}

\textbf{Byte-Pair Encoding.} BPE zostało wprowadzone jako algorytm
kompresji danych przez Gage~(1994), a~adaptowane do tokenizacji
w~NLP przez Sennrich et al.~(2016). Stało się standardem
w~modelach takich jak GPT-2, GPT-3 czy BART. BPE skaluje się
iteracyjnie, łącząc najczęstsze pary znaków aż do osiągnięcia
docelowego rozmiaru słownika.

\textbf{Morfologicznie informowana tokenizacja.} Kłopot z~BPE
dla języków o~bogatej fleksji (polski, czeski, węgierski, turecki)
polega na tym, że scalania oparte na częstotliwości ignorują
strukturę morfologiczną --- tokeny mogą przekraczać granice
morfemów, tworząc jednostki lingwistycznie nienaturalne.
Kilka prac podjęło ten problem:

\begin{itemize}
  \item \textbf{MorphBPE} (Asgari et al., 2025) --- wprowadza
        morfologicznie świadome rozszerzenie BPE z~karami
        za przekraczanie granic morfemów. Propozycje dwóch metryk
        ewaluacji: Morphological Consistency F1 oraz Morphological
        Edit Distance --- tych samych metryk używa nasza praca.
        Eksperymenty na angielskim, rosyjskim, węgierskim
        i~arabskim z~modelami 300M i~1B parametrów.
        \emph{Język polski nie jest testowany.}
  \item \textbf{MorphPiece} (Jabbar, 2023) --- lingwistyczny
        tokenizator oparty na tabeli morfemów z~wewnętrznym BPE
        jako fallback; zaprojektowany dla angielskiego.
  \item \textbf{SentencePiece + Unigram} (Kudo, 2018) --- alternatywa
        dla BPE, modeluje prawdopodobieństwo tokenów zamiast
        częstotliwości par; lepiej radzi sobie z~rzadkimi morfemami.
  \item \textbf{Morfessor} (Creutz \& Lagus, 2007) --- klasyczny
        system morfologicznej segmentacji, ale nie oparty na BPE.
\end{itemize}

\textbf{Polska tokenizacja NLP.} Morfeusz2 (Woliński, 2014)
jest standardowym analizatorem morfologicznym dla języka polskiego.
Wiącek et al.~(2024) zaproponowali benchmark NLPre-PL
dla polskich narzędzi NLP, ale bez ewaluacji morfologicznej
tokenizatorów sub-word. Dawidek~(2026) w~serii preprintów
analizuje ograniczenia BPE w~polskim: wprowadza metryki
stem integrity (SI), inflectional suffix separation (ISS)
i~mean fragmentation per lexeme (MFL) dla 109 form fleksyjnych
10 leksemów. Jej praca pokazuje, że tokenizatory trenowane
na polskim (HerBERT BPE-PL, MFL$=1{,}22$) istotnie przewyższają
tokenizatory angielskie (GPT-2 BPE-EN, MFL$=3{,}91$),
ale \emph{nie testuje morfologicznie informowanego BPE}
jako alternatywy.

\textbf{Wkład naszej pracy.} Asgari et al.~(2025) zaproponowali
MorphBPE i~MorphED, ale nie ewaluowali na polskim. Dawidek~(2026)
analizuje jakość polskich tokenizatorów, ale nie testuje
morfologicznie prowadzonego BPE. Nasza praca łączy oba wątki:
stosuje MorphBPE z~granicami sufiksów z~Morfeusz2 na polskim
korpusie, ewaluuje intrinsic (MorphED z~bootstrap CI)
i~extrinsic z~dwiema architekturami (LSTM i~Transformer)
przy dwóch rozmiarach słownika, oraz przeprowadza analizę
błędów per-suffix-group identyfikującą klasy morfemów,
które zyskują i~tracą na morfologicznym ograniczeniu.
W~najlepszej naszej wiedzy jest to pierwsze zastosowanie
MorphBPE z~analizatorem Morfeusz2 dla języka polskiego
z~ewaluacją extrinsic i~analizą błędów per-suffix.

% ---------------------------------------------------------------
\section{Przegląd zbadanych podejść}
\label{sec:exploration}

Wybór MorphBPE (\texttt{strict\_suffix}) jako metody ewaluowanej
w~tej pracy nie był arbitralny. Przed finalizacją eksperymentów
zbadaliśmy dziesięć podejść do tokenizacji, obejmujących trzy
paradygmaty: częstotliwościowy, semantyczny i~morfologiczny.
Poniżej krótko charakteryzujemy każde podejście i~uzasadniamy
kryteria selekcji.

\textbf{1. Standard BPE} (Sennrich et al., 2016) --- baseline,
klasyczny algorytm BPE oparty na częstotliwości par znaków.
Punkt odniesienia dla wszystkich pozostałych wariantów.

\textbf{2. Parallel BPE} --- nasz wariant eksperymentalny:
równoległe rozważanie $k$ kandydatów na scalenie w~każdej
iteracji (zamiast jednego najczęstszego), z~wyborem po
fertility gain. Cel: sprawdzenie, czy zachłanne wybieranie
najczęstszej pary jest optymalne. Wynik: brak istotnej
poprawy nad standard BPE przy znacznie większym koszcie
obliczeniowym.

\textbf{3. Semantic BPE} --- nasz wariant eksperymentalny:
BPE sterowany osadzeniami word2vec (Mikolov et al., 2013).
Scalenia preferują pary o~wysokim podobieństwie kosinusowym
między osadzeniami części. Testowano $\alpha \in \{0{,}1,
0{,}25, 0{,}5, 1{,}0\}$. Wynik: poprawa alignmentu semantycznego,
ale brak poprawy cross-entropy LM; osadzenia word2vec nie
przechwytują struktury fleksyjnej polskiego.

\textbf{4. Morpheme-aware BPE} --- nasz wczesny prototyp:
segmentacja morfemowa Morfeusz2 jako krok pre-tokenizacji
przed BPE. Wynik: nadmierna fragmentacja --- Morfeusz2 dzieli
słowa na jednostki mniejsze niż użyteczne dla BPE, zwiększając
fertility bez korzyści dla LM.

\textbf{5. MorphBPE variants} (Asgari et al., 2025,
arXiv:2502.00894) --- morfologicznie informowany BPE
z~konfigurowalną siłą ograniczenia (strict, soft 75/50/25\%)
i~typem granic (prefix, suffix, both). Testowano 7 wariantów.
Najlepszy alignment przy zachowaniu fertility: \texttt{strict\_suffix}
($\alpha = 1{,}0$, pełne blokowanie).

\textbf{6. Bayesian BPE} --- nasz wariant eksperymentalny:
scoring bayesowski z~priorytetami lingwistycznymi (polskie
prefiksy i~sufiksy jako prior). Wynik: poprawa alignmentu
dla słów z~prefiksami, ale pogorszenie dla słów bez prefiksu;
mniej spójny niż MorphBPE.

\textbf{7. Compression-aware BPE} (Limisiewicz et al., 2026,
arXiv:2605.01188) --- BPE optymalizujący compression rate
(bytes/token) zamiast częstotliwości par, inspirowany pracą
FAIR/Meta o~compute-optimal tokenization. Wynik: lepsza
kompresja, ale ignoruje granice morfemów --- nie rozwiązuje
problemu alignmentu morfologicznego.

\textbf{8. MorphTok / CBPE} (Brahma et al., 2025,
arXiv:2504.10335) --- morfologiczna pre-tokenizacja
+ Constrained BPE (CBPE) zaprojektowane dla języków indyjskich
(sandhi splitting, dependent vowels). Zaadaptowano na polski
z~Morfeusz2 zamiast sandhi. Wynik: CBPE redukuje fertility,
ale mechanizm zaprojektowany dla abugidy nie przenosi się
dobrze na alfabet łaciński.

\textbf{9. MorphBPE + ToaST} (Asgari et al., 2025 +
Schmidt et al., 2026, arXiv:2605.22705) --- MorphBPE
zastąpiony przez ToaST (Tokenization with Split Trees):
top-down recursive splitting z~byte n-gram counts. Wynik:
ToaST osiąga lepszą kompresję, ale nie ma wbudowanego
morfologicznego ograniczenia; łączenie z~MorphBPE nie
przyniosło synergii.

\textbf{10. LangMAP} (Salhan, Szablewski, Lesci, Buttery,
Meister, Pimentel, 2026, arXiv:2606.23566) --- language-adaptive
MAP tokenization: rozszerzenie UnigramLM o~per-language
distributions. Wynik: poprawa alignmentu morfologicznego,
ale wymaga per-language training; nie jest wariantem BPE
i~nie pozwala na bezpośrednie porównanie z~MorphBPE.

\medskip
\noindent
\textbf{Kryteria selekcji.} MorphBPE \texttt{strict\_suffix}
($\alpha = 1{,}0$) został wybrany, ponieważ:
(i)~utrzymuje fertility na poziomie standardowego BPE
(w~odróżnieniu od Morpheme-aware BPE i~MorphTok),
(ii)~respektuje granice sufiksów z~Morfeusz2 (w~odróżnieniu
od Semantic BPE, Bayesian BPE i~Compression-aware BPE),
(iii)~jest prosty i~kompatybilny ze standardowym inference BPE
(w~odróżnieniu od ToaST i~LangMAP),
(iv)~pochodzi z~peer-reviewed paperu (ACL 2025 Findings)
z~ewaluacją na 4 językach, ale \emph{bez polskiego} --- co
uzasadnia nasz wkład.

\begin{table}[h]
  \centering
  \caption{Faza eksploracyjna: fertility badanych podejść
  (korpus 5\,000 przemówień, vocab=500). Fertility jest jedyną
  metryką porównywalną między prototypami (MorphED liczony różnymi
  implementacjami w~poszczególnych skryptach). Baseline Standard BPE
  $=2{,}45$. Niżej = bardziej kompaktowa tokenizacja.}
  \label{tab:exploration}
  \small
  \begin{tabular}{lcl}
    \toprule
    Podejście & Fertility & Werdykt \\
    \midrule
    \rowcolor{slayergreen!20}
    Standard BPE (baseline)      & 2{,}45 & punkt odniesienia \\
    Semantic BPE ($\alpha{=}0{,}5$) & 2{,}46 & brak poprawy CE LM \\
    Morpheme-aware BPE           & 2{,}54 & nadmierna fragmentacja \\
    MorphBPE \texttt{strict\_prefix} & 2{,}53 & pogarsza alignment \\
    MorphBPE \texttt{strict\_both}   & 2{,}54 & gorszy niż suffix \\
    \rowcolor{slayergreen!20}
    MorphBPE \texttt{strict\_suffix} & 2{,}46 & \textbf{wybrany} \\
    LangMAP                      & 3{,}06 & wymaga per-language, nie-BPE \\
    ToaST                        & 3{,}24 & brak ograniczenia morf. \\
    MorphTok / CBPE              & 3{,}58 & fragmentacja (abugida) \\
    \bottomrule
  \end{tabular}
\end{table}

Tab.~\ref{tab:exploration} podsumowuje mierzalne wyniki fazy
eksploracyjnej. Tylko \texttt{strict\_suffix} osiąga fertility
na poziomie baseline'u BPE ($2{,}46$ vs $2{,}45$) przy jednoczesnym
respektowaniu granic morfologicznych; podejścia zwiększające
fertility (LangMAP, ToaST, MorphTok) wprowadzają nadmierną
fragmentację. Warianty \texttt{Parallel BPE}, \texttt{Bayesian BPE}
i~\texttt{Compression-aware BPE} oceniono jakościowo, bez zapisanych
metryk ilościowych.

% ---------------------------------------------------------------
\section{Wprowadzenie}

\textbf{Pytanie:} czy wprowadzenie morfologicznych ograniczeń
do algorytmu BPE poprawia jakość tokenizacji dla języka polskiego
--- zarówno w~sensie alignmentu morfologicznego (intrinsic),
jak i~użyteczności dla modelu języka (extrinsic)?

\smallskip
\noindent
Nitka dialektyczna:

\medskip
\noindent
\textbf{Teza~(T):} MorphBPE (\texttt{strict\_suffix}) produkuje tokeny
lepiej zgrane z~granicami morfemów niż standardowy BPE, co przekłada
się na niższą cross-entropy modelu języka LSTM przy tym samym
rozmiarze słownika.

\medskip
\noindent
\textbf{Antyteza~(A):} Morfologiczne ograniczenia hamują częstotliwościowe
scalanie BPE, zwiększając fertility (liczbę tokenów na słowo)
i~wzmacniając fragmentację, co może pogorszyć -- nie poprawić --
cross-entropy LM.

\begin{slayerinsight}
\textbf{Synteza~(S):} MorphBPE suffix utrzymuje fertility na poziomie
BPE ($\Delta < 0{,}1\%$) przy obu rozmiarach słownika, jednocześnie
redukując MorphED o~3{,}6\% (vocab=500) i~6{,}8\% (vocab=1000).
Cross-entropy LSTM spada o~1{,}3\% i~2{,}6\% odpowiednio;
Transformer potwierdza trend ($-1{,}3\%$ i~$-2{,}7\%$).
Morfologiczne ograniczenia nie kosztują compression, a~pomagają LM.
\end{slayerinsight}

% ---------------------------------------------------------------
\section{Metoda}

\subsection*{Korpus}
Korpus: 10\,000 przemówień z~posiedzeń Sejm RP (polski parlament),
pochodzących z~publicznego API sejm.gov.pl, udostępnionych
poprzez repozytorium HuggingFace\footnote{\url{https://huggingface.co/datasets/PiotrSty/sejm-speeches-corpus}}.
Podział train/test: 80/20 (8000/2000), seed=42.
Tekst: małe litery, tokenizacja whitespace (interpunkcja
pozostaje przyklejona do słów).

\subsection*{Tokenizatory}
\textbf{Standard BPE:} klasyczny algorytm BPE oparty na częstotliwości
par. Słownik inicjowany pojedynczymi znakami; scalanie iteracyjne
najczęstszych par aż do osiągnięcia docelowego rozmiaru słownika.

\smallskip
\noindent
\textbf{MorphBPE (\texttt{strict\_suffix}):} wariant BPE z~morfologicznym
ograniczeniem. Granice sufiksów pozyskiwane z~Morfeusz2.
Dla każdej pary tokenów w~słowie obliczana jest kara:
jeśli scalenie przekracza granicę sufiksu, jego częstotliwość
mnożona jest przez $(1 - \alpha)$. W~wariancie \emph{strict}
($\alpha = 1{,}0$, używanym w~głównych eksperymentach) oznacza to
pełne zablokowanie scaleń przekraczających granicę sufiksu.
Warianty miękkie ($\alpha < 1$) oraz over-penalization ($\alpha > 1$)
badamy systematycznie w~sekcji~\ref{sec:penalty-sweep}.
Ograniczenie preferuje scalenia wewnątrz rdzeni i~sufiksów,
zniechęcając scalenia przekraczające granice (cross-boundary merges).

\subsection*{Przykłady tokenizacji}
Tab.~\ref{tab:examples} pokazuje konkretne słowa i~jak je tnie
BPE vs MorphBPE suffix (vocab=1000). Różnice są subtelne --- oba
tokenizatory produkują podobne segmenty dla większości słów.
Kluczowa różnica widać przy \emph{samorządowych}:
BPE tnie na \texttt{sa~|~morządowych} (przekracza granicę
rdzeń$\to$sufiks), MorphBPE tnie na \texttt{sa~|~morządowy~|~ch}
(respektuje granicę sufiksu \texttt{-ch}).
Dla większości słów oba tokenizatory produkują identyczne segmentacje
--- różnica ujawnia się tylko tam, gdzie BPE „przeskakuje''
granicę morfemu.

\begin{table}[h]
  \centering
  \caption{Przykłady tokenizacji: BPE vs MorphBPE suffix (vocab=1000).
  Pionowa kreska \texttt{|} oznacza granicę tokena.}
  \label{tab:examples}
  \small
  \begin{tabular}{lll}
    \toprule
    Słowo & BPE & MorphBPE suffix \\
    \midrule
    przedstawienie   & przedstawie~|~nie       & przedstawie~|~nie \\
    rozważania       & rozważania              & rozważania \\
    odpowiedzialności & odpowiedzial~|~ności    & odpowiedzial~|~ności \\
    najważniejszych  & naj~|~waż~|~niej~|~szych & naj~|~waż~|~niej~|~szych \\
    współpracy       & współ~|~pracy           & współ~|~pracy \\
    niezależności    & nie~|~zale~|~ż~|~ności   & nie~|~zale~|~ż~|~ności \\
    przekształcenia  & prze~|~ksz~|~tał~|~cenia & prze~|~ksz~|~tał~|~cenia \\
    \rowcolor{slayergray}
    samorządowych    & sa~|~morządowych        & sa~|~morządowy~|~ch \\
    \bottomrule
  \end{tabular}
\end{table}

Wyróżniony wiersz (\emph{samorządowych}) ilustruje kluczową różnicę:
BPE łączy rdzeń z~sufiksem w~jeden token (\texttt{morządowych}),
podczas gdy MorphBPE zachowuje granicę sufiksu \texttt{-ch},
dzieląc na \texttt{morządowy~|~ch}. To pozwala LM rozpoznać
wzorzec sufiksu \texttt{-ch} wspólny dla wielu przypadków
(genitiv/accusativ plural) i~lepiej generalizować.

\subsection*{Rozmiary słownika}
Eksperymenty przeprowadzono dla dwóch rozmiarów słownika:
$V = 500$ i~$V = 1000$.
Rozmiar słownika obejmuje znaki bazowe + scalone tokeny.

\subsection*{Ewaluacja intrinsic}
\textbf{Fertility} (Fert): liczba tokenów na słowo.
Niżej = bardziej kompaktowa tokenizacja.
\[
  \mathrm{Fert} = \frac{|\mathrm{tokens}|}{|\mathrm{words}|}
\]

\textbf{Morphological Edit Distance} (MorphED): średnia odległość
Levenshteina między segmentacją tokenizatora a~segmentacją
morfologiczną (Morfeusz2) po normalizacji.
Niżej = lepszy alignment morfologiczny.
\[
  \mathrm{MorphED} = \frac{1}{|W|} \sum_{w \in W}
    d_{\mathrm{Lev}}\!\bigl(\mathrm{tok}(w),\; \mathrm{morph}(w)\bigr)
\]

\subsection*{Ewaluacja extrinsic}
\textbf{Model języka LSTM:} jednowarstwowy LSTM
($\mathrm{embed\_dim}=64$, $\mathrm{hidden\_dim}=128$),
trenowany na tokenach z~danego tokenizatora.
Parametry: 3 epoki, batch=128, seq\_len=32, lr=0.003, optimizer Adam.
Dane treningowe: 500 tekstów (train), test: 200 tekstów (test).

\smallskip
\noindent
\textbf{Model języka Transformer:} 2-warstwowy Transformer encoder
($\mathrm{embed\_dim}=64$, 4 głowice, $\mathrm{ff\_dim}=256$),
trenowany na tych samych danych. Parametry: 3 epoki, batch=128,
seq\_len=32, lr=0.001, optimizer AdamW, dropout=0.1.
Positional embeddings uczone wspólnie z~modelowaniem.

\textbf{Cross-entropy} (CE): średnia entropia krzyżowa na token
na zbiorze testowym.
\textbf{Perplexity} (PPL): $= \exp(\mathrm{CE})$.
Niżej = lepszy model języka.

\subsection*{Istotność statystyczna}
Bootstrap confidence intervals: 1000 resamples, 95\% CI.
Resampling po słowach ze zwracaniem z~zbioru testowego.
Raportowane wartości: średnia $[\mathrm{CI}_{2{,}5\%},\, \mathrm{CI}_{97{,}5\%}]$.

% ---------------------------------------------------------------
\section{Wyniki}
\label{sec:results}

\subsection*{Ewaluacja intrinsic}

\begin{table}[h]
  \centering
  \caption{Ewaluacja intrinsic: Fertility i~MorphED z~bootstrap CI (95\%).
  Korpus: 10\,000 przemówień sejmowych, split 80/20, seed=42.
  Bielik-11B podany jako zewnętrzny punkt odniesienia (produkcyjny tokenizator).}
  \label{tab:intrinsic}
  \begin{tabular}{lrccr}
    \toprule
    Metoda & Vocab & Fertility [CI] & MorphED [CI] & Czas [s] \\
    \midrule
    BPE                  & 500  & 2{,}443 [2{,}439, 2{,}447] & 2{,}503 [2{,}497, 2{,}509] & 368{,}9 \\
    MorphBPE suffix      & 500  & 2{,}444 [2{,}440, 2{,}449] & 2{,}414 [2{,}408, 2{,}420] & 469{,}5 \\
    \addlinespace
    BPE                  & 1000 & 1{,}932 [1{,}929, 1{,}936] & 1{,}838 [1{,}832, 1{,}843] & 654{,}9 \\
    MorphBPE suffix      & 1000 & 1{,}934 [1{,}930, 1{,}937] & 1{,}712 [1{,}707, 1{,}717] & 1031{,}8 \\
    \addlinespace
    BPE ($\alpha{=}0$)   & 8\,000 & 1{,}131 & 0{,}574 [0{,}569, 0{,}578] & --- \\
    MorphBPE ($\alpha{=}2$) & 8\,000 & 1{,}149 & \textbf{0{,}456} [0{,}451, 0{,}460] & --- \\
    \addlinespace
    BPE ($\alpha{=}0$)   & 16\,000 & 1{,}051 & 0{,}442 [0{,}438, 0{,}446] & --- \\
    MorphBPE ($\alpha{=}2$) & 16\,000 & 1{,}077 & \textbf{0{,}338} [0{,}335, 0{,}342] & --- \\
    \addlinespace
    BPE ($\alpha{=}0$)   & 32\,000 & 1{,}020 & 0{,}389 [0{,}385, 0{,}393] & --- \\
    MorphBPE ($\alpha{=}2$) & 32\,000 & 1{,}049 & \textbf{0{,}308} [0{,}304, 0{,}311] & --- \\
    \addlinespace
    \rowcolor{slayergray}
    Bielik-11B orig. (ref.) & 32\,128 & 3{,}220 & --- & --- \\
    \rowcolor{slayergray}
    Bielik APT4 (ref.)   & $\sim$32k & 1{,}620 & --- & --- \\
    \bottomrule
  \end{tabular}
\end{table}

\textbf{Fertility.} MorphBPE suffix utrzymuje fertility na poziomie BPE:
$\Delta = +0{,}04\%$ (vocab=500) i~$\Delta = +0{,}10\%$ (vocab=1000).
Różnice wewnątrz CI --- brak istotnego kosztu compression.

\textbf{Punkt odniesienia zewnętrzny.}
Dla kontekstualizacji wyników dołączamy Bielik-11B-v2 --- produkcyjny
tokenizator (SentencePiece BPE, vocab=32\,128) wytrenowany na wielomiliardowym
korpusie polskim. Przy porównywalnym rozmiarze słownika ($\sim$32k)
MorphBPE osiąga fertility $3{,}1\times$ niższą niż oryginalny Bielik
(1{,}049 vs 3{,}220) i~$1{,}5\times$ niższą niż polsko-zoptymalizowany
APT4 (1{,}049 vs 1{,}620). Oryginalny tokenizer Bielika (Mistral-derived,
vocab=32\,128) jest używany w~Bielik-11B v2, v2.x i~v3.
APT4 to nowy, polsko-zoptymalizowany tokenizer opisany w~raporcie
o~optymalizacji tokenizera Bielik v3.
Bezpośrednie porównanie MorphED nie jest możliwe (Bielik nie
publikuje segmentacji morfologicznej), ale różnica fertility wskazuje,
że MorphBPE produkuje znacznie bardziej kompaktową tokenizację.
Uwaga: Bielik trenowany jest na innym, większym korpusie --- porównanie
jest orientacyjne i~służy wyłącznie jako punkt odniesienia skali.

\textbf{MorphED.} MorphBPE suffix redukuje MorphED istotnie:
\begin{itemize}
  \item vocab=500: $2{,}414$ vs $2{,}503$ ($-3{,}6\%$), CI nie pokrywają się.
  \item vocab=1000: $1{,}712$ vs $1{,}838$ ($-6{,}8\%$), CI nie pokrywają się.
  \item vocab=8\,000: $0{,}456$ vs $0{,}574$ ($-20{,}5\%$), CI nie pokrywają się.
  \item vocab=16\,000: $0{,}338$ vs $0{,}442$ ($-23{,}4\%$), CI nie pokrywają się.
  \item vocab=32\,000: $0{,}308$ vs $0{,}389$ ($-21{,}0\%$), CI nie pokrywają się.
\end{itemize}
Efekt rośnie z~rozmiarem słownika --- przy większym vocab BPE tworzy
dłuższe tokeny przekraczające granice morfemów, podczas gdy MorphBPE
utrzymuje granice sufiksów. Redukcja MorphED jest 3--6$\times$ silniejsza
przy vocab=8k--32k niż przy 500--1000.

\begin{figure}[h]
  \centering
  \begin{tikzpicture}
    \begin{axis}[
        ybar,
        bar width=14pt,
        width=0.82\linewidth, height=5.5cm,
        symbolic x coords={vocab=500, vocab=1000},
        xtick=data,
        ymin=1.5, ymax=2.7,
        ylabel={MorphED (niżej = lepiej)},
        nodes near coords,
        every node near coord/.style={font=\small},
        xmajorgrids=true, ymajorgrids=true,
        grid style={gray!20},
        legend style={at={(0.5,-0.28)}, anchor=north, legend columns=2, font=\small},
        enlarge x limits=0.5,
    ]
      \addplot[fill=slayerblue!60, draw=slayerblue] coordinates {
        (vocab=500, 2.503)
        (vocab=1000, 1.838)
      };
      \addlegendentry{BPE}
      \addplot[fill=slayergreen!60, draw=slayergreen] coordinates {
        (vocab=500, 2.414)
        (vocab=1000, 1.712)
      };
      \addlegendentry{MorphBPE suffix}
    \end{axis}
  \end{tikzpicture}
  \caption{MorphED --- BPE vs MorphBPE suffix. CI 95\% nie pokrywają się
  przy obu rozmiarach słownika.}
  \label{fig:morphed}
\end{figure}

\subsection*{Ewaluacja extrinsic}

\begin{table}[h]
  \centering
  \caption{Ewaluacja extrinsic: cross-entropy i~perplexity LSTM LM.
  Model: LSTM(64, 128, 1 warstwa), 3 epoki, batch=128, seq\_len=32.}
  \label{tab:extrinsic}
  \begin{tabular}{lrccr}
    \toprule
    Metoda & Vocab & Cross-Entropy & Perplexity & Czas LM [s] \\
    \midrule
    BPE                  & 500  & 3{,}2758 & 26{,}46 & 698{,}2 \\
    MorphBPE suffix      & 500  & 3{,}2332 & 25{,}36 & 794{,}7 \\
    \addlinespace
    BPE                  & 1000 & 3{,}7438 & 42{,}26 & 810{,}5 \\
    MorphBPE suffix      & 1000 & 3{,}6461 & 38{,}33 & 789{,}5 \\
    \bottomrule
  \end{tabular}
\end{table}

\textbf{Cross-entropy.} MorphBPE suffix obniża CE względem BPE:
\begin{itemize}
  \item vocab=500: $3{,}2332$ vs $3{,}2758$ ($-1{,}3\%$).
  \item vocab=1000: $3{,}6461$ vs $3{,}7438$ ($-2{,}6\%$).
\end{itemize}

\textbf{Perplexity.} MorphBPE suffix obniża PPL:
\begin{itemize}
  \item vocab=500: $25{,}36$ vs $26{,}46$ ($-4{,}2\%$).
  \item vocab=1000: $38{,}33$ vs $42{,}26$ ($-9{,}3\%$).
\end{itemize}

\begin{figure}[h]
  \centering
  \begin{tikzpicture}
    \begin{axis}[
        ybar,
        bar width=14pt,
        width=0.82\linewidth, height=5.5cm,
        symbolic x coords={vocab=500, vocab=1000},
        xtick=data,
        ymin=3.0, ymax=4.0,
        ylabel={Cross-entropy (niżej = lepiej)},
        nodes near coords,
        every node near coord/.style={font=\small},
        xmajorgrids=true, ymajorgrids=true,
        grid style={gray!20},
        legend style={at={(0.5,-0.28)}, anchor=north, legend columns=2, font=\small},
        enlarge x limits=0.5,
    ]
      \addplot[fill=slayerblue!60, draw=slayerblue] coordinates {
        (vocab=500, 3.2758)
        (vocab=1000, 3.7438)
      };
      \addlegendentry{BPE}
      \addplot[fill=slayergreen!60, draw=slayergreen] coordinates {
        (vocab=500, 3.2332)
        (vocab=1000, 3.6461)
      };
      \addlegendentry{MorphBPE suffix}
    \end{axis}
  \end{tikzpicture}
  \caption{Cross-entropy LSTM LM --- BPE vs MorphBPE suffix.
  MorphBPE suffix obniża CE o~1{,}3\% (vocab=500) i~2{,}6\% (vocab=1000).}
  \label{fig:ce}
\end{figure}

\textbf{Uwaga:} przy vocab=1000 oba tokenizatory osiągają wyższą CE
niż przy vocab=500. Wynika to z~większej liczby rzadkich tokenów
przy ograniczonym rozmiarze danych treningowych LM (500 tekstów)
i~małej pojemności modelu. Efekt ten jest spójny
dla obu metod --- różnica względna MorphBPE vs BPE rośnie jednak
z~rozmiarem słownika ($-1{,}3\% \to -2{,}6\%$ dla LSTM,
$-1{,}3\% \to -2{,}7\%$ dla Transformera), co wskazuje,
że korzyść z~morfologicznego ograniczenia jest silniejsza
przy większym słowniku.

\subsection*{Ewaluacja Transformer LM}

\begin{table}[h]
  \centering
  \caption{Ewaluacja extrinsic: cross-entropy i~perplexity Transformer LM.
  MorphBPE suffix: wariant strict ($\alpha = 1{,}0$).
  Model: Transformer(64, 4h, 2L, ff=256), 3 epoki, batch=128, seq\_len=32.}
  \label{tab:transformer}
  \begin{tabular}{lrccr}
    \toprule
    Metoda & Vocab & Cross-Entropy & Perplexity & Czas LM [s] \\
    \midrule
    BPE                  & 500  & 3{,}1306 & 22{,}89 & 1093{,}6 \\
    MorphBPE suffix      & 500  & 3{,}0902 & 21{,}98 & 1457{,}9 \\
    \addlinespace
    BPE                  & 1000 & 3{,}5250 & 33{,}95 & 962{,}9 \\
    MorphBPE suffix      & 1000 & 3{,}4306 & 30{,}89 & 1361{,}6 \\
    \bottomrule
  \end{tabular}
\end{table}

\textbf{Cross-entropy (Transformer).} MorphBPE suffix obniża CE:
\begin{itemize}
  \item vocab=500: $3{,}0902$ vs $3{,}1306$ ($-1{,}3\%$).
  \item vocab=1000: $3{,}4306$ vs $3{,}5250$ ($-2{,}7\%$).
\end{itemize}
Trend jest spójny z~wynikami LSTM: MorphBPE obniża CE przy obu
rozmiarach słownika, a~korzyść rośnie z~$V$. Magnituda jest
zbliżona do LSTM ($-1{,}3\%$ przy vocab=500, $-2{,}7\%$ przy vocab=1000),
co potwierdza, że korzyść z~morfologicznego ograniczenia jest
odporna na wybór architektury (LSTM vs Transformer z~attention).

\subsection*{Analiza błędów per-suffix}

Aby zrozumieć, \emph{które} klasy morfemów zyskują i~tracą
na morfologicznym ograniczeniu, przeprowadziliśmy analizę
per-suffix-group cross-entropy na modelu Transformer (vocab=1000).
Dla każdej grupy słów kończących się danym sufiksem obliczyliśmy
średni per-token CE osobno dla BPE i~MorphBPE.

\begin{table}[h]
  \centering
  \caption{Analiza błędów: średni per-token CE wg grupy sufiksu
  (Transformer, vocab=1000, $\alpha = 0{,}5$ --- exploracyjna analiza
  z~miękkim ograniczeniem dla lepszego rozróżnienia grup).
  $\Delta < 0$ = MorphBPE lepszy.}
  \label{tab:errors}
  \small
  \begin{tabular}{lrccr}
    \toprule
    Sufiks & N & BPE CE & MorphBPE CE & $\Delta$ [\%] \\
    \midrule
    \rowcolor{slayergreen!20}
    -ach    & 49 & 4{,}3432 & 3{,}8443 & $-11{,}5$ \\
    \rowcolor{slayergreen!20}
    -ego    & 38 & 4{,}6068 & 4{,}4511 & $-3{,}4$ \\
    \rowcolor{slayergreen!20}
    -ości   & 47 & 5{,}1136 & 4{,}9541 & $-3{,}1$ \\
    \rowcolor{slayergreen!20}
    -enia   & 26 & 2{,}6377 & 2{,}6033 & $-1{,}3$ \\
    -owania & 31 & 1{,}6421 & 1{,}6443 & $+0{,}1$ \\
    -ami    & 50 & 3{,}8532 & 3{,}8720 & $+0{,}5$ \\
    \rowcolor{slayerred!20}
    -ny     & 43 & 3{,}0789 & 3{,}3573 & $+9{,}0$ \\
    \rowcolor{slayerred!20}
    -ów     & 32 & 2{,}7220 & 3{,}0345 & $+11{,}5$ \\
    \rowcolor{slayerred!20}
    -ych    & 40 & 1{,}9407 & 2{,}7585 & $+42{,}1$ \\
    \rowcolor{slayerred!20}
    -owych  & 37 & 1{,}9336 & 3{,}0116 & $+55{,}8$ \\
    \rowcolor{slayerred!20}
    -ć      & 49 & 3{,}5789 & 5{,}5765 & $+55{,}8$ \\
    \bottomrule
  \end{tabular}
\end{table}

\textbf{Wzorzec:} MorphBPE zyskuje na sufiksach złożonych
(wieloznakowych: \texttt{-ach}, \texttt{-ego}, \texttt{-ości},
\texttt{-enia}), gdzie respektowanie granicy morfemu pozwala LM
lepiej generalizować wzorce fleksyjne. MorphBPE traci na sufiksach
krótkich (\texttt{-ć}, \texttt{-ych}, \texttt{-ów}), gdzie wymuszone
oddzielenie jednoliterowego lub dwuliterowego sufiksu prowadzi
do nieefektywnej fragmentacji --- LM musi modelować krótki token
w~izolacji, bez kontekstu rdzenia.

To sugeruje, że \emph{adaptacyjna kara} --- silniejsza dla sufiksów
długich, słabsza dla krótkich --- mogłaby połączyć korzyści
morfologicznego alignmentu z~mniejszą fragmentacją.

Odwołania: rys.~\ref{fig:hero} (hero MorphED),
rys.~\ref{fig:morphed} (MorphED),
rys.~\ref{fig:ce} (cross-entropy LSTM),
tab.~\ref{tab:intrinsic}, tab.~\ref{tab:extrinsic},
tab.~\ref{tab:transformer}, tab.~\ref{tab:errors}.

% ---------------------------------------------------------------
\section{Wpływ siły kary i~typu ograniczenia}
\label{sec:penalty-sweep}

Aby zbadać, jak siła kary $\alpha$ i~typ ograniczenia
(prefix, suffix, both) wpływają na jakość tokenizacji
i~wydajność modelu języka, przeprowadziliśmy systematyczny
przegląd $\alpha \in \{0{,}0;\, 0{,}1;\, 0{,}25;\, 0{,}5;\, 0{,}75;\,
1{,}0\}$ $\times$ 3 typy ograniczenia $\times$ 2 rozmiary słownika
= 36 eksperymentów (z~których $\alpha = 0{,}0$ jest równoważne
standardowemu BPE niezależnie od typu ograniczenia).
Następnie rozszerzyliśmy zakres o~$\alpha > 1{,}0$ (tzw.\
over-penalization): $\alpha \in \{1{,}0;\, 1{,}5;\, 2{,}0;\, 2{,}5;\,
3{,}0\}$ dla ograniczenia suffix $\times$ 2 vocab = 10 dodatkowych
eksperymentów. Wszystkie eksperymenty używają tego samego korpusu,
podziału train/test i~architektury Transformer LM (64d, 4h, 2L,
ff=256, 3 epoki, 500/200 tekstów).

\subsection*{Wyniki główne ($\alpha \in [0;\,1]$)}

\begin{table}[h]
  \centering
  \caption{Wpływ $\alpha$ i~typu ograniczenia na cross-entropy
  Transformer LM (vocab=500). $\alpha = 0{,}0$ = standardowy BPE.}
  \label{tab:sweep500}
  \small
  \begin{tabular}{lcccccc}
    \toprule
    Constraint & $\alpha{=}0{,}0$ & $\alpha{=}0{,}1$ & $\alpha{=}0{,}25$ & $\alpha{=}0{,}5$ & $\alpha{=}0{,}75$ & $\alpha{=}1{,}0$ \\
    \midrule
    suffix & 3{,}131 & 3{,}144 & 3{,}142 & 3{,}111 & 3{,}095 & \textbf{3{,}090} \\
    prefix & 3{,}131 & 3{,}138 & 3{,}140 & 3{,}138 & 3{,}139 & 3{,}143 \\
    both   & 3{,}131 & 3{,}144 & 3{,}132 & 3{,}135 & 3{,}116 & 3{,}111 \\
    \bottomrule
  \end{tabular}
\end{table}

\begin{table}[h]
  \centering
  \caption{Wpływ $\alpha$ i~typu ograniczenia na cross-entropy
  Transformer LM (vocab=1000).}
  \label{tab:sweep1000}
  \small
  \begin{tabular}{lcccccc}
    \toprule
    Constraint & $\alpha{=}0{,}0$ & $\alpha{=}0{,}1$ & $\alpha{=}0{,}25$ & $\alpha{=}0{,}5$ & $\alpha{=}0{,}75$ & $\alpha{=}1{,}0$ \\
    \midrule
    suffix & 3{,}525 & 3{,}545 & 3{,}523 & 3{,}477 & 3{,}462 & \textbf{3{,}431} \\
    prefix & 3{,}525 & 3{,}522 & 3{,}522 & 3{,}513 & 3{,}517 & 3{,}511 \\
    both   & 3{,}525 & 3{,}520 & 3{,}505 & 3{,}501 & 3{,}469 & 3{,}437 \\
    \bottomrule
  \end{tabular}
\end{table}

\textbf{Trzy wnioski z~zakresu $\alpha \le 1$:}

\begin{enumerate}
  \item \textbf{Suffix >> prefix >> both.} Ograniczenie suffix
        poprawia CE monotonicznie ze wzrostem $\alpha$
        ($-1{,}3\%$ przy vocab=500, $-2{,}7\%$ przy vocab=1000
        dla $\alpha = 1{,}0$). Ograniczenie prefix \emph{pogarsza}
        CE względem BPE przy vocab=500 i~daje marginalną poprawę
        przy vocab=1000. Typ \texttt{both} jest gorszy od
        \texttt{suffix} alone --- komponent prefiksowy przeszkadza.
        To potwierdza typologię polskiego jako języka sufiksalnego:
        fleksja realizowana jest na końcu wyrazu, prefiksy są rzadkie.
  \item \textbf{Niska kara ($\alpha = 0{,}1$) szkodzi.} Przy obu
        rozmiarach słownika $\alpha = 0{,}1$ daje \emph{gorsze} CE
        niż BPE ($+0{,}4\%$ i~$+0{,}6\%$). Słaba kara zaburza
        naturalną kolejność scalenia BPE bez wystarczająco silnego
        sygnału morfologicznego --- najgorszy z~obu światów.
  \item \textbf{Korzyść rośnie z~$\alpha$ i~z~$V$.} Im silniejsza
        kara i~im większy słownik, tym większa poprawa CE.
        Przy $\alpha = 1{,}0$ i~vocab=1000: $\Delta\mathrm{CE} = -0{,}094$
        ($-2{,}7\%$), podczas gdy przy vocab=500: $\Delta = -0{,}041$
        ($-1{,}3\%$).
\end{enumerate}

\subsection*{Over-penalization ($\alpha > 1{,}0$)}

W~oryginalnej formulacji MorphBPE wartość $\alpha = 1{,}0$ oznacza
pełne blokowanie scaleń przekraczających granice morfemów
(częstotliwość pary mnożona przez $1 - \alpha = 0$). Wartości
$\alpha > 1{,}0$ dają \emph{ujemny} score dla takich par
($\mathrm{freq} \times (1 - \alpha) < 0$). W~praktyce oznacza to,
że gdy \emph{wszystkie} dostępne pary przekraczają granice,
najczęstsza z~nich (least negative) jest wybierana zamiast losowej
--- oraz że crossing merges są aktywnie ``odpychane'' w~kolejności
scalenia, faworyzując merges wewnątrz morfemów.

\begin{table}[h]
  \centering
  \caption{Over-penalization: $\alpha > 1{,}0$, ograniczenie suffix.
  CE i~MorphED z~bootstrap CI (95\%).}
  \label{tab:extended}
  \small
  \begin{tabular}{lrcccc}
    \toprule
    $\alpha$ & Vocab & Fertility & MorphED [CI] & CE & PPL \\
    \midrule
    1{,}0  & 500  & 2{,}449 & 2{,}423 [2{,}419, 2{,}427] & 3{,}092 & 22{,}01 \\
    1{,}5  & 500  & 2{,}443 & 2{,}418 [2{,}413, 2{,}422] & 3{,}092 & 22{,}03 \\
    \rowcolor{slayergreen!20}
    2{,}0  & 500  & 2{,}388 & 2{,}332 [2{,}327, 2{,}336] & \textbf{3{,}044} & \textbf{21{,}00} \\
    2{,}5  & 500  & 2{,}388 & 2{,}330 [2{,}326, 2{,}334] & 3{,}046 & 21{,}04 \\
    3{,}0  & 500  & 2{,}387 & 2{,}322 [2{,}318, 2{,}327] & 3{,}046 & 21{,}02 \\
    \addlinespace
    1{,}0  & 1000 & 1{,}935 & 1{,}717 [1{,}713, 1{,}721] & 3{,}432 & 30{,}93 \\
    1{,}5  & 1000 & 1{,}929 & 1{,}697 [1{,}693, 1{,}701] & 3{,}420 & 30{,}57 \\
    2{,}0  & 1000 & 1{,}891 & 1{,}637 [1{,}633, 1{,}641] & 3{,}368 & 29{,}02 \\
    2{,}5  & 1000 & 1{,}870 & 1{,}599 [1{,}595, 1{,}602] & 3{,}330 & 27{,}94 \\
    \rowcolor{slayergreen!20}
    3{,}0  & 1000 & 1{,}869 & 1{,}592 [1{,}589, 1{,}596] & \textbf{3{,}324} & \textbf{27{,}77} \\
    \bottomrule
  \end{tabular}
\end{table}

\textbf{Kluczowe odkrycie:} $\alpha = 2{,}0$ daje \emph{skokową}
poprawę nad $\alpha = 1{,}0$, nie plateau:

\begin{itemize}
  \item vocab=500: CE $3{,}092 \to 3{,}044$ ($\Delta = -0{,}048$,
        $-1{,}6\%$) --- większa niż cała poprawa od $\alpha = 0{,}5$
        do $\alpha = 1{,}0$ ($\Delta = 0{,}020$).
  \item vocab=1000: CE $3{,}432 \to 3{,}368$ ($\Delta = -0{,}064$,
        $-1{,}9\%$), a~$\alpha = 3{,}0$ daje $3{,}324$
        ($\Delta = -0{,}108$ vs $\alpha = 1{,}0$, $-3{,}1\%$).
  \item MorphED również skacze: $2{,}423 \to 2{,}332$ (vocab=500),
        $1{,}717 \to 1{,}592$ (vocab=1000).
  \item Fertility spada: $2{,}449 \to 2{,}388$ (vocab=500) ---
        over-penalization tworzy bardziej konserwatywne segmentacje.
\end{itemize}

\textbf{Plateau ok.~$\alpha \in [2{,}0;\, 3{,}0]$.} Dalsze zwiększanie
$\alpha$ powyżej 2{,}0 przynosi marginalne zyski (CE $3{,}044 \to
3{,}046$ przy vocab=500), ale przy vocab=1000 trend jest nadal
malejący ($3{,}368 \to 3{,}324$). Optimum wydaje się $\alpha \approx
2{,}0$--$3{,}0$ w~zależności od rozmiaru słownika.

\textbf{Mechanizm.} Przy $\alpha \le 1{,}0$ crossing pairs otrzymują
score $\ge 0$ --- są blokowane tylko wtedy, gdy istnieją alternatywy
z~dodatnim score. Przy $\alpha > 1{,}0$ crossing pairs dostają
score $< 0$, co ma dwa efekty: (1)~gdy wszystkie pary crossing,
najczęstsza jest wybierana (least negative) zamiast losowej;
(2)~kolejność scalenia jest reorganizowana --- merges wewnątrz
morfemów są faworyzowane nawet kosztem opóźnienia crossing merges.
To tłumaczy spadek fertility (bardziej konserwatywne scalenia)
i~poprawę MorphED (lepszy alignment).

\begin{figure}[h]
  \centering
  \begin{tikzpicture}
    \begin{axis}[
        width=0.85\linewidth, height=6cm,
        xlabel={Kara $\alpha$},
        ylabel={Cross-entropy (niżej = lepiej)},
        xmin=-0.1, xmax=3.1,
        ymin=3.0, ymax=3.6,
        xtick={0,0.5,1,1.5,2,2.5,3},
        grid=major, grid style={gray!20},
        legend style={at={(0.98,0.98)}, anchor=north east, font=\small},
    ]
      \addplot[slayerblue, mark=*, thick] coordinates {
        (0,3.131) (0.5,3.111) (1,3.090) (1.5,3.092) (2,3.044) (2.5,3.046) (3,3.046)
      };
      \addlegendentry{vocab=500}
      \addplot[slayerred, mark=square*, thick] coordinates {
        (0,3.525) (0.5,3.477) (1,3.432) (1.5,3.420) (2,3.368) (2.5,3.330) (3,3.324)
      };
      \addlegendentry{vocab=1000}
      \draw[dashed, gray] (axis cs:1,3.0) -- (axis cs:1,3.6);
      \node[font=\small, gray] at (axis cs:1,3.55) {over-pen.};
    \end{axis}
  \end{tikzpicture}
  \caption{Wpływ kary $\alpha$ na cross-entropy Transformer LM.
  Linia przerywana oddziela standardowy zakres ($\alpha \le 1$)
  od over-penalization ($\alpha > 1$). Skok przy $\alpha = 2{,}0$
  jest wyraźnie widoczny przy obu rozmiarach słownika.}
  \label{fig:alpha-sweep}
\end{figure}

\textbf{Implikacja dla tej pracy.} Główne eksperymenty
(sek.~\ref{sec:results}, tab.~\ref{tab:intrinsic}--\ref{tab:transformer})
używają wariantu \emph{strict} ($\alpha = 1{,}0$) --- pełnego blokowania
scaleń przekraczających granicę sufiksu. Wyniki sweepu pokazują jednak,
że $\alpha = 2{,}0$ (over-penalization) jest konsekwentnie lepszy
we wszystkich metrykach. Wartości dla $\alpha = 1{,}0$ pozostają ważne
jako silny punkt odniesienia, ale optymalna konfiguracja to
\texttt{suffix} z~$\alpha = 2{,}0$.

\subsection*{Znaczenie praktyczne}

Co w~praktyce oznaczają te odkrycia?

\textbf{1. Mniej tokenów, tańsze przetwarzanie.}
Over-penalization ($\alpha = 2{,}0$) obniża fertility z~2{,}449 do 2{,}388
tokenów na słowo (vocab=500) --- to około 2{,}5\% mniej tokenów na ten sam tekst.
W~skali infrastruktury LLM, gdzie koszt inference rośnie liniowo
z~liczbą tokenów, 2{,}5\% redukcji tokenów przekłada się
na około 2{,}5\% tańsze generowanie i~krótszy czas odpowiedzi.
Przy vocab=1000 redukcja jest jeszcze większa: 1{,}935 $\to$ 1{,}869
($-3{,}4\%$).

\textbf{2. Lepszy model języka przy tym samym budżecie.}
Cross-entropy spada o~1{,}6\%--3{,}1\% względem $\alpha = 1{,}0$,
co oznacza, że~LM popełnia mniej błędów predykcji bez zwiększenia
liczby parametrów, rozmiaru danych ani mocy obliczeniowej treningu.
To ``darmowa'' poprawa --- jedynym kosztem jest zmiana jednego
hiperparametru ($\alpha$) podczas treningu tokenizatora.

\textbf{3. Korzyść rośnie ze słownikiem.}
Przy vocab=1000 poprawa CE nad $\alpha = 1{,}0$ wynosi $-3{,}1\%$
($\alpha = 3{,}0$), podczas gdy przy vocab=500 tylko $-1{,}6\%$.
W~praktyce produkcyjne tokenizatory używają słowników 30\,000--100\,000
tokenów. Jeśli trend utrzyma się, korzyść z~over-penalization może być
tam jeszcze większa --- ale wymaga to weryfikacji eksperymentalnej.

\textbf{4. Ekonomia tokenizacji.}
Tokenizator jest trenowany raz, a~używany miliardy razy. Koszt treningu
MorphBPE z~$\alpha = 2{,}0$ jest identyczny jak z~$\alpha = 0{,}5$
--- to ten sam algorytm, różny jest tylko mnożnik w~funkcji score.
Zatem: zerowy koszt dodatkowy przy treningu, stała korzyść
przy każdym inference. W~kontekście modeli obsługujących języki
o~bogatej fleksji (polski, czeski, węgierski, turecki, fiński),
optymalizacja $\alpha$ to niskowiszący owoc --- prosta zmiana
dająca mierzalną poprawę jakości i~kosztów.

\textbf{5. Typologia języka ma znaczenie.}
Wyniki sweepu constraint type potwierdzują, że~dla języków
sufiksalnych (jak polski) kluczowe jest respektowanie granic sufiksów,
a~nie prefiksów. To ma implikację projektową: dla danego języka
należy wybrać constraint type zgodny z~jego typologią morfologiczną,
a~nie stosować uniwersalne ``both''.

% ---------------------------------------------------------------
\section{Skalowanie do rozmiarów produkcyjnych}
\label{sec:scaling}

Eksperymenty główne (sek.~\ref{sec:intrinsic}, \ref{sec:extrinsic})
używały słowników 500 i~1000 tokenów --- wystarczających do pokazania
trendu, ale dalekich od rozmiarów produkcyjnych (30\,000--100\,000).
Aby zweryfikować, czy korzyść z~MorphBPE utrzymuje się przy większych
słownikach, przeprowadziliśmy skalowanie do vocab $\in \{8000, 16\,000, 32\,000\}$
przy $\alpha \in \{0{,}0;\, 1{,}0;\, 2{,}0\}$, na pełnym zbiorze treningowym
(8000 tekstów, backend heurystyczny).

\begin{table}[h]
\centering
\small
\begin{tabular}{rrrrrrr}
\hline
Vocab & $\alpha$ & ActVocab & Fertility & MorphED & CI (95\%) & Czas \\
\hline
8000 & 0{,}0 & 8000 & 1{,}131 & 0{,}574 & [0{,}569, 0{,}578] & 2{,}3\,h \\
8000 & 1{,}0 & 8000 & 1{,}141 & 0{,}504 & [0{,}499, 0{,}509] & 1{,}8\,h \\
8000 & 2{,}0 & 8000 & 1{,}149 & \textbf{0{,}456} & [0{,}451, 0{,}460] & 1{,}6\,h \\
16000 & 0{,}0 & 16000 & 1{,}051 & 0{,}442 & [0{,}438, 0{,}446] & 2{,}9\,h \\
16000 & 1{,}0 & 16000 & 1{,}064 & 0{,}387 & [0{,}384, 0{,}391] & 3{,}2\,h \\
16000 & 2{,}0 & 16000 & 1{,}077 & \textbf{0{,}338} & [0{,}335, 0{,}342] & 3{,}2\,h \\
32000 & 0{,}0 & 32000 & 1{,}020 & 0{,}389 & [0{,}385, 0{,}393] & 6{,}0\,h \\
32000 & 1{,}0 & 32000 & 1{,}035 & 0{,}347 & [0{,}343, 0{,}351] & 5{,}3\,h \\
32000 & 2{,}0 & 32000 & 1{,}049 & \textbf{0{,}308} & [0{,}304, 0{,}311] & 5{,}9\,h \\
\hline
\end{tabular}
\caption{Skalowanie MorphBPE suffix do rozmiarów produkcyjnych.
Train: 8000 tekstów, test: 500 tekstów, backend heurystyczny.}
\label{tab:scaling}
\end{table}

\textbf{Wyniki.}
Przy wszystkich rozmiarach słownika MorphBPE redukuje MorphED
względem BPE ($\alpha = 0$) o~20{,}5\%--23{,}4\%, przy czym
CI 95\% nie pokrywają się przy żadnym rozmiarze.
Over-penalization ($\alpha = 2{,}0$) daje dodatkowe 9{,}5\%--12{,}6\%
redukcji MorphED względem $\alpha = 1{,}0$.
Najlepszy wynik w~całym badaniu: vocab=32\,000, $\alpha = 2{,}0$,
MorphED$=0{,}308$.

\textbf{Koszt fertility.}
Przy vocab=500--1000 koszt fertility wynosi $\Delta < 0{,}1\%$.
Przy skalowaniu koszt rośnie: $+1{,}6\%$ (8k), $+2{,}5\%$ (16k), $+2{,}8\%$ (32k).
Nadal jest to niewielki koszt, ale trend wskazuje, że~przy bardzo dużych
słownikach ($> 100\,000$) może stać się istotny.

\textbf{Porównanie z~Bielik-11B.}
Oryginalny tokenizer Bielika (Mistral-derived, vocab=32\,128), używany
w~Bielik-11B v2, v2.x i~v3, osiąga fertility$=3{,}22$ na~reprezentatywnym
tekście polskim. Nowy, polsko-zoptymalizowany tokenizer APT4
($\sim$32k vocab) obniża ten wskaźnik do $1{,}62$.
MorphBPE przy vocab=32\,000 osiąga fertility$=1{,}05$ ---
$3{,}1\times$ niższą niż oryginał i~$1{,}5\times$ niższą niż APT4.
Bezpośrednie porównanie MorphED nie jest możliwe (Bielik nie publikuje
segmentacji morfologicznej), ale różnica fertility wskazuje, że~MorphBPE
produkuje znacznie bardziej kompaktową tokenizację przy tym samym
rozmiarze słownika.

\textbf{Wniosek.}
Korzyść z~MorphBPE nie tylko utrzymuje się przy skalowaniu ---
jest \emph{silniejsza} niż przy małych słownikach.
Redukcja MorphED rośnie z~$-3{,}6\%$ (vocab=500) do $-21{,}0\%$ (vocab=32\,000).
Over-penalization ($\alpha = 2{,}0$) pozostaje optymalną konfiguracją
na wszystkich rozmiarach.

% ---------------------------------------------------------------
\section{Warunek obalenia}

Teza o~wyższości MorphBPE suffix nad BPE upada, jeśli spełniony
zostanie którykolwiek z~warunków:

\begin{enumerate}
  \item MorphBPE suffix zwiększa fertility o~$> 1\%$ względem BPE
        przy którymś rozmiarze słownika.
        Aktualnie: $\Delta = +0{,}04\%$ (500), $+0{,}10\%$ (1000) --- \textbf{nie spełniony}.
  \item MorphBPE suffix zwiększa MorphED względem BPE
        (CI się pokrywają lub MorphBPE > BPE).
        Aktualnie: CI nie pokrywają się, MorphBPE $<$ BPE --- \textbf{nie spełniony}.
  \item MorphBPE suffix zwiększa cross-entropy LM względem BPE
        (LSTM lub Transformer).
        Aktualnie: CE MorphBPE $<$ CE BPE przy obu vocab i~obu architekturach
        --- \textbf{nie spełniony}.
\end{enumerate}

Żaden z~warunków nie jest spełniony --- synteza nie zostaje obalona.

% ---------------------------------------------------------------
\section{Dyskusja}

\textbf{Główny wynik.} MorphBPE suffix dominuje nad standardowym BPE
na wszystkich mierzonych metrykach (MorphED, cross-entropy, perplexity)
przy obu rozmiarach słownika i~obu architekturach LM, bez istotnego
kosztu fertility. Korzyść rośnie z~rozmiarem słownika:
MorphED $-3{,}6\% \to -6{,}8\%$,
CE (LSTM) $-1{,}3\% \to -2{,}6\%$,
CE (Transformer) $-1{,}3\% \to -2{,}7\%$,
PPL (LSTM) $-4{,}2\% \to -9{,}3\%$.
PPL (Transformer) $-4{,}0\% \to -9{,}0\%$.
Skalowanie do vocab=32\,000 (sek.~\ref{sec:scaling}) potwierdza
trend: redukcja MorphED rośnie do $-21{,}0\%$, a~over-penalization
pozostaje optymalne na wszystkich rozmiarach.
Przegląd siły kary (sek.~\ref{sec:penalty-sweep}) pokazuje ponadto,
że optymalna wartość to \texttt{suffix} z~$\alpha = 2{,}0$
(over-penalization), która obniża CE o~kolejne $1{,}6\%$--$3{,}1\%$
względem sztywnego blokowania ($\alpha = 1{,}0$).

\textbf{Dlaczego działa.} Standardowy BPE optymalizuje wyłącznie
częstotliwość par, co prowadzi do scaleń przekraczających granice
morfemów (np.\ \texttt{prze+staw} zamiast \texttt{prze+stawienie}).
MorphBPE suffix kształtuje przestrzeń scaleń tak, aby granice
sufiksów były respektowane. Tokeny wynikowe są bardziej zgrane
z~podziałem morfologicznym, co ułatwia LM generalizację
(podobne sufiksy $\to$ podobne reprezentacje).

\textbf{Ograniczenia.}
(1)~Ewaluacja extrinsic używa małych modeli (LSTM 128, Transformer 64/2L)
i~ograniczonego zbioru treningowego (500 tekstów). Wynik może się różnić
dla większych modeli i~pełnego zbioru (8000 tekstów).
(2)~MorphBPE wymaga analizatora morfologicznego (Morfeusz2) ---
dodatkowa zależność, niedostępna dla wszystkich języków.
(3)~Analiza per-suffix pokazuje, że MorphBPE traci na sufiksach krótkich
(\texttt{-ć}, \texttt{-ych}); adaptacyjna kara może ten problem złagodzić.
(4)~Over-penalization ($\alpha > 1$) została przetestowana tylko
na ograniczeniu suffix; prefix i~both z~$\alpha > 1$ pozostają do zbadania.
(5)~Wszystkie eksperymenty sweepu używają Transformer LM o~tej samej
architekturze; ewaluacja z~LSTM i~większymi modelami może dać inne optimum $\alpha$.

\textbf{Co dalej.}
(1)~Ewaluacja extrinsic (LM) przy vocab=32\,000 --- skalowanie intrinsic
(sek.~\ref{sec:scaling}) potwierdza korzyść, ale ewaluacja z~modelem
języka przy produkcyjnych rozmiarach słownika wymaga większych modeli
(Transformer 4+ warstw) i~pełnego zbioru treningowego.
(2)~Adaptacyjna kara $\alpha$ zależna od długości sufiksu
(silniejsza dla długich, słabsza dla krótkich) --- analiza
per-suffix (tab.~\ref{tab:errors}) sugeruje, że krótkie sufiksy
(\texttt{-ć}, \texttt{-ych}) tracą na sztywnym ograniczeniu.
(3)~Test na innych językach słowiańskich (czeski, rosyjski)
z~odpowiednimi analizatorami morfologicznymi.
(4)~Porównanie z~SentencePiece i~Unigram LM tokenizer.
(5)~Over-penalization ($\alpha > 1$) dla typów \texttt{prefix}
i~\texttt{both} oraz weryfikacja, czy optymum $\alpha \approx 2{,}0$
przenosi się na większe modele i~pełny zbiór danych.

% ---------------------------------------------------------------
\section*{Podpis}
Dokument autorstwa: \emph{Piotr Styła}.
Korpus: Polish Sejm Corpus (sejm.gov.pl, licencja publiczna).
Kod: \texttt{paper\_experiments.py}, \texttt{lm\_eval.py},
\texttt{transformer\_eval.py}, \texttt{penalty\_sweep.py},
\texttt{penalty\_sweep\_extended.py}.
Skompilowany w~CI (\texttt{.tex}~$\to$~PDF).
Data: \today.
